\chapter{CΩSMOS: The Physical Universe as RMG}
\label{sec:cosmos-rmg}

This is the chapter where we commit a mild act of cosmological treason.

So far we’ve:

- built an ontology (RMG + DPO),
- given it geometry (MRMW, curvature, worldlines),
- showed how physics-like behavior \emph{emerges} from this machinery,
- and compiled ordinary physical theories into RMG universes.

That’s all cute.

Now we point the gun at the only universe we actually have and say:

\begin{quote}
\textbf{There exists some RMG $G_\ast$ and some rule set $\mathcal{R}_\ast$ such that the pair $(G_\ast, \mathcal{R}_\ast)$ is\emd{}up to coarse-graining\emd{}our cosmos.}
\end{quote}

We’ll call this hypothetical pair \emph{CΩSMOS}. This chapter is not "here is the final rule of the universe." It’s:

- what CΩSMOS would have to look like to be taken seriously, and
- how you’d go about hunting it inside MRMW.

It’s a spec, not a solution.

\subsubsection*{27.1  What Does It Mean to "Be Our Universe"?}

We need a precise target.

Saying "encode the universe" is useless; we need constraints an RMG can pass or fail.

Let’s define:

\begin{quote}
\textbf{A CΩSMOS candidate} is a pair $(G,\mathcal{R})$ plus a coarse-graining $C$ such that, for some choice of units and error tolerances, the induced macro-dynamics reproduce the empirical regularities we call "physics".
\end{quote}

Concretely, that means:

\begin{enumerate}
  \item \textbf{Spacetime-like behavior.} There exists an emerging manifold-ish structure with:
    \begin{itemize}
      \item approximate $3+1$-dimensionality,
      \item a Lorentzian causal structure (light cones, relativity of simultaneity),
      \item and metric behavior matching GR at accessible scales.
    \end{itemize}

  \item \textbf{Matter + interactions.} There exist substructures whose coarse behavior matches:
    \begin{itemize}
      \item observed particle spectra (masses, charges, spins),
      \item observed gauge interactions (at least an SM-like gauge structure),
      \item and known coupling strengths (within experimental precision).
    \end{itemize}

  \item \textbf{Quantum statistics.} Rewrite bundles and their weighting reproduce:
    \begin{itemize}
      \item interference patterns,
      \item Bell-type correlations,
      \item and the usual Born-rule frequencies for measurement outcomes.
    \end{itemize}

  \item \textbf{Cosmological history.} There exist initial conditions such that:
    \begin{itemize}
      \item coarse-grained early-time states look like a hot, dense phase,
      \item expansion, structure formation, and background radiation emerge with the right patterns,
      \item and large-scale homogeneity / isotropy are approximately respected.
    \end{itemize}
\end{enumerate}

If your RMG doesn’t tick boxes like this, it’s not "our universe." It’s just fanfic.

\subsubsection*{27.2  Design Constraints from Observation}

Before we sketch CΩSMOS, let’s list the non-negotiables as design requirements on $(G,\mathcal{R})$.

\paragraph{(1) Locality with bounded influence.}

We need:

- DPO rules whose matches are bounded-radius patterns,
- a provable (or at least strong) upper bound on causal influence speed,
- so we can recover something light-cone-like.

Formally: there exists a $v_\text{max}$ such that the causal cone around any rewrite event grows no faster than linearly with "time" (rewrite depth).

\paragraph{(2) Approximate Lorentz invariance.}

The causal graph of rewrite events must admit:

- a family of coarse-grained frames,
- in which light-cone structure is invariant under pseudo-orthogonal transformations,
- and in which no global preferred frame is detectable at accessible scales.

In graph terms: large-scale statistics of causal neighborhoods should be highly symmetric, with anisotropies confined to small-scale lattice artifacts or extreme regimes.

\paragraph{(3) Discrete but effectively continuous.}

At fundamental level, CΩSMOS is discrete (finite-degree graphs, finite information per local pattern). But:

- there must exist large-scale coarse-grainings that look smooth,
- with effective field equations that approximate known PDEs,
- and with discretization artifacts below current experimental bounds.

\paragraph{(4) Gauge and internal symmetries.}

Automorphism groups of local graph motifs must contain something isomorphic to:

- a product of continuous gauge groups (or a good discrete approximation),
- acting on "matter" subgraphs in the right representations.

Symmetries cannot be tacked on externally; they must be actual RMG symmetries / rule equivariances.

\paragraph{(5) A natural low-entropy "initial" region.}

We need an explanation for:

- a simple, low-entropy early configuration,
- that expands and becomes more complex under $\mathcal{R}$,
- without fine-tuning so extreme it’s indistinguishable from magic.

In other words: the Big Bang is a special region of MRMW, not a special pleading.

These are obnoxiously strong constraints. Good. That’s how you filter most of MRMW out.

\subsubsection*{27.3  Pre-Geometry: Causal Graph First, Coordinates Later}

A sane move for CΩSMOS is to go \emph{pre-geometric}:

- Start with a causal RMG: a graph whose primary structure is "who can affect whom," not coordinates.
- Let spacetime, distance, and dimension fall out as emergent properties of that causal structure.

So we posit:

\begin{itemize}
  \item A base graph $G_\text{cause}$:
    \begin{itemize}
      \item nodes = events (rewrite applications),
      \item edges = causal influence (as defined in earlier chapters),
      \item optional higher-order structure for "simultaneity slices" or foliations.
    \end{itemize}

  \item A rule set $\mathcal{R}_\text{pre}$ that:
    \begin{itemize}
      \item generates new events from local configurations of prior events,
      \item respects bounded influence,
      \item preserves some global invariants (for energy-, momentum-like quantities).
    \end{itemize}
\end{itemize}

Actual "geometry’’ then is derived:

- Dimension is estimated from volume-growth vs. radius in the causal graph.
- Metric distances are defined by geodesic lengths or appropriate functionals on paths.
- Curvature is defined by deviations of volume/angle relations from flat expectations.

You don’t start with $\mathbb{R}^{3,1}$. You discover something that \emph{acts} like it.

RMG recursion gives you a clean layering:

- Level 0: raw events and their causal edges.
- Level 1: coarse-grained "regions" and emergent coordinates.
- Level 2: fields and matter living on these emergent structures.

CΩSMOS is this whole tower, not just the bottom layer.

\subsubsection*{27.4  From Causal Graph to Spacetime}

We now sketch the map:

\[
  (G_\text{cause}, \mathcal{R}_\text{pre}) \quad\longrightarrow\quad \text{effective spacetime}.
\]

Given a large causal graph:

\begin{enumerate}
  \item \textbf{Extract slices.} Find maximal antichains (sets of events mutually spacelike) as candidate "equal-time" surfaces.

  \item \textbf{Define adjacency.} Build a derived graph whose nodes are small regions of these slices; connect them if they share many causal links to the same future/past events.

  \item \textbf{Estimate dimension.} Measure how the number of nodes within $k$ steps grows with $k$; fit an effective dimension $d$ from scaling laws.

  \item \textbf{Fit a metric.} Treat the causal structure as giving you a conformal class of metrics; use additional combinatorial data (like density of events) to fix scale.

  \item \textbf{Read curvature.} Quantify deviations from flat scaling; map them onto curvature tensors in the emergent manifold picture.
\end{enumerate}

If CΩSMOS is viable, this procedure should spit out:

- $d \approx 3$ spatial + 1 "time" dimension,
- nearly flat large-scale behavior,
- localized curvature where matter/energy-like structures concentrate.

GR’s Einstein equation then becomes a \emph{phenomenological fit} to how density of certain graph structures correlates with deviations in causal volume growth.

We’re not deriving Einstein’s equation here; we’re marking the target: any $(G_\ast,\mathcal{R}_\ast)$ that wants to be CΩSMOS must make something like Einstein inevitable in its continuum limit.

\subsubsection*{27.5  Matter and Fields as Internal Structure}

With spacetime-like geometry extracted, we decorate it.

In CΩSMOS:

\begin{itemize}
  \item \textbf{Matter} is represented by persistent, localized subgraphs that:
    \begin{itemize}
      \item propagate along emergent geodesics,
      \item interact via local rewrite at intersections,
      \item carry conserved charges encoded as structural invariants.
    \end{itemize}

  \item \textbf{Gauge fields} are represented as:
    \begin{itemize}
      \item labels on edges of the emergent adjacency graph,
      \item or small subgraphs attached to adjacency motifs,
      \item with group-like composition around cycles (plaquettes).
    \end{itemize}

  \item \textbf{Internal symmetries} show up as automorphisms of these attached structures that leave the causal skeleton invariant.
\end{itemize}

The Standard Model then corresponds to the claim:

\begin{quote}
Near our universe’s region in MRMW, the local automorphism structure of matter+field subgraphs looks like an SM-like gauge group acting on a zoo of persistent patterns we call "particles".
\end{quote}

You don’t hardcode $SU(3)\times SU(2)\times U(1)$; you search for rule sets whose graph symmetries \emph{are} that (or something observationally equivalent).

\subsubsection*{27.6  Gravity as Curvature in Causal Rewrite}

We already defined curvature of MRMW; now we specialize:

- \textbf{Spacetime curvature} is curvature of the emergent causal manifold.
- \textbf{Gravitational interaction} is the effect of that curvature on matter subgraph worldlines.

In CΩSMOS there are (at least) two plausible ways this can arise:

\paragraph{(A) Structural coupling.}

Rules in $\mathcal{R}_\ast$ directly couple matter density to the pattern of future causal connections:

- where certain patterns ("mass-energy") concentrate,
- the local statistics of new event creation change (more/less branching, altered connectivity),
- leading to modified geodesic behavior in the emergent metric.

Einstein-like equations are then statistical summaries of this coupling.

\paragraph{(B) Entropic / combinatorial gravity.}

Gravity emerges as:

- a tendency of rewrite trajectories to explore regions of MRMW with maximal microstate volume under macroscopic constraints,
- leading to effective attraction between matter subgraphs because those configurations dominate combinatorially.

You can mix these pictures; both live happily in RMG.

The critical bit: you never add "gravity" as a separate force. You let curvature of the causal structure \emph{be} what we call gravity, with matter patterns as the source of that curvature.

\subsubsection*{27.7  Constants and Scales as Combinatorial Parameters}

Fundamental constants become less mystical and more annoying:

- The \textbf{speed of light} $c$ is the maximum effective information propagation rate imposed by bounded-influence rules; numerically, it’s a ratio of spacelike to timelike scaling in the causal graph.

- \textbf{Coupling constants} are combinatorial parameters controlling:
  \begin{itemize}
    \item relative frequencies of different local rewrite motifs,
    \item weights/amplitudes assigned to particular rewrite bundles,
    \item or structural biases in how new edges/nodes are created.
  \end{itemize}

- \textbf{Mass scales} show up as:
  \begin{itemize}
    \item effective inertia of certain subgraphs (how readily their patterns change),
    \item or phases in their rewrite amplitudes (how their histories interfere).
  \end{itemize}

Planck units are just the natural units of CΩSMOS’s combinatorics. Once you’ve fixed:

- the fundamental discrete step size (graph scale),
- the basic "tick" of rewrite (event depth),
- and the amplitude weighting scheme,

all the other constants are dimensionless knobs in rule space.

The "why these values?" question then becomes:

\begin{quote}
Why does our universe’s rule set $\mathcal{R}_\ast$ sit at this particular point in rulial space, instead of nearby ones where those combinatorial knobs are different?
\end{quote}

That’s not answered here, but at least now it’s clearly a \emph{ruliometric} question, not metaphysical noise.

\subsubsection*{27.8  Quantum CΩSMOS: Bundles All the Way Down}

To make CΩSMOS actually resemble our universe, you must inject the Part~III machinery:

- \textbf{Rewrite bundles} represent superposition: multiple compatible histories of $(G,\mathcal{R})$ evolve in parallel in MRMW.

- \textbf{Amplitudes} are associated with histories via an action functional:
  \[
    \mathcal{A}(\text{history}) \propto e^{i S[\text{history}]/\hbar},
  \]
  where $S$ is computed from local graph motifs along the path.

- \textbf{Interference} arises when different histories lead to the same coarse-grained RMG state; their amplitudes sum (constructively or destructively) before probabilities are taken.

- \textbf{Measurement} is implemented as:
  \begin{itemize}
    \item a special class of observer subgraphs that entangle their internal state with branch structure,
    \item followed by effective "collapse" via minimal-path selection or decoherence in the observer’s own coarse-graining.
  \end{itemize}

In CΩSMOS, then, there is no ontological distinction between "classical region" and "quantum region." There are only:

- regimes where bundles are thick and interference is visible,
- regimes where coarse-graining has effectively washed interference out.

Our universe looks "quasi-classical" on macroscales because the observer-subgraphs we inhabit are violently coarse-grained; not because the underlying RMG forgot how to do quantum.

\subsubsection*{27.9  Where Are We in MRMW?}

If CΩSMOS is one point in MRMW, we can ask: what’s its neighborhood?

Nearby in rulial space you’d find:

\begin{itemize}
  \item universes with slightly different constants (same qualitative structure),
  \item universes with the same causal pre-geometry but different matter subgraphs,
  \item universes with same low-energy behavior but wildly different high-energy rules,
  \item and universes that are "dual" to ours: different microscopic rules, but equivalent macroscopic physics under aggressive coarse-graining.
\end{itemize}

This gives a concrete sense to:

- \textbf{Anthropic questions}: are we in a rare region where observers are even possible?
- \textbf{Dualities}: are different theoretical descriptions just different coordinates on the same region of MRMW?
- \textbf{Fine-tuning}: how thin is the sliver of rulial space that produces a CΩSMOS-like universe?

In principle, CΩMPUTER can treat these as engineering problems:

- sample rule sets near a candidate $(G_\ast,\mathcal{R}_\ast)$,
- measure resulting coarse behavior,
- map out the "anthropically habitable" basin in MRMW.

That’s not philosophy; that’s a rulial Monte Carlo.

\subsubsection*{27.10  How CΩSMOS Becomes a Research Program}

Right now, all of this is a clean story and an ugly TODO.

Turning CΩSMOS from concept to science means:

\begin{enumerate}
  \item \textbf{Specifying families of candidate rule sets} $\mathcal{R}_\theta$ with tunable parameters $\theta$ (topology, local motifs, weighting schemes).

  \item \textbf{Evolving them at scale} in a CΩMPUTER runtime that:
    \begin{itemize}
      \item natively handles RMG + DPO + bundles,
      \item tracks causal structure and curvature,
      \item supports aggressive coarse-graining and observable queries.
    \end{itemize}

  \item \textbf{Defining a fitness function} that scores $(G,\mathcal{R})$ on:
    \begin{itemize}
      \item dimension and causal structure,
      \item symmetry and conservation properties,
      \item spectrum and interaction patterns of matter-like structures,
      \item cosmological large-scale behavior.
    \end{itemize}

  \item \textbf{Searching MRMW} with:
    \begin{itemize}
      \item guided exploration (gradient-like methods in rule space),
      \item adversarial universes (Part~IV) to stress-test candidate laws,
      \item and rulial provenance (Part~IV) to track why a candidate works.
    \end{itemize}
\end{enumerate}

In other words: "find the rule of the universe’’ becomes a large-scale inverse problem on MRMW, with CΩMPUTER as the experimental apparatus.

\subsubsection*{27.11  Why Put Physics Inside CΩMPUTER at All?}

Because once you do, several walls fall at once:

\begin{itemize}
  \item The wall between simulation and reality: there’s just one substrate (RMG + DPO), and "real" vs "modeled" are about which region of MRMW you’re in and how you’re coupled to it.

  \item The wall between physical and non-physical systems: markets, organisms, AIs, and galaxies all become different RMG universes with different constraints.

  \item The wall between "fundamental" and "effective" theories: both are just different coarse-grainings / coordinates on the same underlying combinatorics.
\end{itemize}

CΩSMOS is not a separate product; it’s just the name we give to the particular RMG that earned the right to be called "our universe."

\medskip

The next chapter, \emph{The Future of Intelligence in MRMW}, zooms back in from the cosmos to the agents inside it.

We’ll take everything we’ve said about CΩSMOS and ask:

\begin{quote}
Given that we live inside one particular RMG universe with all this structure, what does it really mean to be "intelligent" here\emd{}and how far can we push that, once CΩMPUTER lets us see and steer our position in MRMW?  
\end{quote}
